% ===== ElegantBook (cyan theme) customizations for 《庖丁解牛 pigo》 =====
% The cyan theme supplies the accent color (\structurecolor). Everything below
% is layered on top: Chinese font fallback, code-block styling, and the
% \begin{experimentbox} environment consumed by experiment_box.lua.

% ── Chinese font fallback across platforms ───────────────────────
% macOS ships Songti SC / Heiti SC; most Linux/CI images ship Noto CJK.
% Degrade gracefully so the build never hard-fails on a missing family.
% Order: (1) display names — resolved by a full TeX Live xelatex and on Linux;
% (2) explicit macOS system-font paths — tectonic can't resolve the CJK display
%     names via fontconfig, so it needs the filename-stem + Path + Extension form;
% (3) Noto CJK for Linux/CI; (4) let xeCJK keep its default as a last resort.
\IfFontExistsTF{Songti SC}{%
  \setCJKmainfont[AutoFakeBold=2.5]{Songti SC}%
  \setCJKsansfont[AutoFakeBold=2.5]{Heiti SC}%
  \setCJKmonofont[AutoFakeBold=2.5]{Songti SC}%
}{\IfFontExistsTF{/System/Library/Fonts/Supplemental/Songti.ttc}{%
  \setCJKmainfont{Songti}[Path=/System/Library/Fonts/Supplemental/,Extension=.ttc,UprightFeatures={FontIndex=0},AutoFakeBold=2.5]%
  \setCJKsansfont{STHeiti Medium}[Path=/System/Library/Fonts/,Extension=.ttc,UprightFeatures={FontIndex=0},AutoFakeBold=2.5]%
  \setCJKmonofont{Songti}[Path=/System/Library/Fonts/Supplemental/,Extension=.ttc,UprightFeatures={FontIndex=0},AutoFakeBold=2.5]%
}{\IfFontExistsTF{Noto Serif CJK SC}{%
  \setCJKmainfont[AutoFakeBold=2.5]{Noto Serif CJK SC}%
  \setCJKsansfont[AutoFakeBold=2.5]{Noto Sans CJK SC}%
  \setCJKmonofont[AutoFakeBold=2.5]{Noto Sans CJK SC}%
}{\IfFontExistsTF{Noto Sans CJK SC}{%
  \setCJKmainfont[AutoFakeBold=2.5]{Noto Sans CJK SC}%
  \setCJKsansfont[AutoFakeBold=2.5]{Noto Sans CJK SC}%
  \setCJKmonofont[AutoFakeBold=2.5]{Noto Sans CJK SC}%
}{%
  % Last resort: let ctex/xeCJK keep its own default so compilation proceeds.
}}}}

% ── Unicode glyph fallback for symbols the Latin text font lacks ──
\usepackage{newunicodechar}
\IfFontExistsTF{Arial Unicode MS}{\newfontfamily\fallbackfont{Arial Unicode MS}}{%
  \IfFontExistsTF{Noto Sans CJK SC}{\newfontfamily\fallbackfont{Noto Sans CJK SC}}{%
    \newfontfamily\fallbackfont{DejaVu Sans}}}
\newunicodechar{✓}{{\fallbackfont ✓}}
\newunicodechar{✗}{{\fallbackfont ✗}}
\newunicodechar{★}{{\fallbackfont ★}}
\newunicodechar{≈}{{\fallbackfont ≈}}
\newunicodechar{→}{{\fallbackfont →}}

% ── Monospace font for code (falls back gracefully) ──────────────
\IfFontExistsTF{Menlo}{\setmonofont[Scale=0.86]{Menlo}}{%
  \IfFontExistsTF{DejaVu Sans Mono}{\setmonofont[Scale=0.86]{DejaVu Sans Mono}}{}}

% ── Accent color: derive chrome shades from the theme's cyan ─────
\colorlet{accent}{structurecolor}
\colorlet{accentbg}{structurecolor!8}
\colorlet{accentrule}{structurecolor!40}

% Clickable internal cross-reference command used by crossref.lua.
\newcommand{\crossreflink}[2]{\hyperref[#1]{\textcolor{structurecolor}{#2}}}

% ── tcolorbox libraries (ElegantBook already loads tcolorbox) ────
\tcbuselibrary{skins,breakable}

% ── Wrap long lines in code blocks so they don't overflow the page ─
\usepackage{fvextra}
\fvset{breaklines=true,breakanywhere=true,%
  breaksymbolleft={\footnotesize\textcolor{accentrule}{$\hookrightarrow$}\,}}
\DefineVerbatimEnvironment{Highlighting}{Verbatim}%
  {commandchars=\\\{\},breaklines,breakanywhere}

% ── Code blocks: subtle framed box (shared by highlighted & plain) ─
\newtcolorbox{codebox}{%
  breakable, enhanced,
  colback=black!3, colframe=accentrule,
  boxrule=0.5pt, leftrule=2.5pt,
  left=0.7em, right=0.6em, top=0.4em, bottom=0.4em,
  arc=0.6mm, before skip=0.7em, after skip=0.7em}
\renewenvironment{Shaded}{\begin{codebox}}{\end{codebox}}

% ── Experiment callout box (rendered by experiment_box.lua) ──────
% Paragraphs/divs marked「实验 N-M」are wrapped in this bordered box.
\newtcolorbox{experimentbox}{%
  breakable, enhanced,
  colback=accentbg, colframe=accent,
  boxrule=0.6pt, left=0.8em, right=0.8em, top=0.6em, bottom=0.6em,
  arc=1mm, before skip=0.9em, after skip=0.9em}
